% Standalone report. Compile from the repository root with:
% pdflatex -interaction=nonstopmode -halt-on-error -output-directory=BUILD doc/synthesis/population/education_model.tex
% Run twice to resolve the table of contents and cross-references.
\documentclass[11pt,a4paper]{article}
\usepackage[T1]{fontenc}
\usepackage[utf8]{inputenc}
\usepackage{lmodern}
\usepackage[margin=23mm]{geometry}
\usepackage{amsmath,booktabs,array,longtable}
\usepackage{listings}
\usepackage[colorlinks=true,linkcolor=blue,urlcolor=blue]{hyperref}
\urlstyle{tt}
\newcommand{\code}[1]{\nolinkurl{#1}}
\newcolumntype{L}[1]{>{\raggedright\arraybackslash}p{#1}}
\setlength{\parindent}{0pt}
\setlength{\parskip}{5pt}
\setlength{\emergencystretch}{3em}
\lstset{basicstyle=\ttfamily\small,breaklines=true,columns=fullflexible,
  keepspaces=true,showstringspaces=false,frame=single}
\title{Education modelling in the Swiss synthetic population\\
\large Implemented methodology and configuration reference}
\author{eqasim-switzerland}
\date{21 October 2026}

\begin{document}
\maketitle

\textbf{Implementation snapshot.} Branch
\code{feature/education-destination-calibration}, commit \code{ed19479}.
This report describes the current \code{census: statpop} workflow, including
the education, apprenticeship and projection changes committed on this date.
It distinguishes active behaviour from diagnostic calculations and earlier
approaches that are no longer used. Defaults below are read from the source
code; individual YAML configurations can override them. No model or population
was changed to prepare this report.

\textbf{Central distinction.} Student status, a usual education location and
attendance on the simulated day are three different outcomes. A student has an
education anchor even if their selected diary contains no education visit.
BFS enrolment provides relative tertiary institution/campus shares. Those
shares are normalised to the modelled cohort; they do not independently fix
the national number of students or the number attending each day.

\tableofcontents
\newpage

\section{Data sources and their roles}

\begin{longtable}{L{.22\textwidth}L{.27\textwidth}L{.43\textwidth}}
\toprule Source & Coverage used & Role and interpretation \\
\midrule\endhead
STATPOP & 2023 person and household extracts & Population identities, ages,
households and registered home coordinates. Selected geocoded people missing
from the private-household table are retained in technical one-person households. \\
Structural survey (SS) & Pooled 2021--2023 records with their survey weights;
2023 is the default stock-audit year & Employment/student models, study
classification, school departure and destination flows, and secondary commute
distance distributions. The education models use age-15+ evidence. \\
Microcensus (MZ) & Existing MZ 2015 pipeline inputs; interviews from age 6 &
Activity-chain donors, diary attendance, usual commute predictors and child
distance references. MZ 2021 comparisons made during investigation are not an
additional active training input in this implementation. \\
BFS enrolment & Academic year 2023/24 by default & Pooled UH and FH/PH institution
totals, including disjoint continuing/other and legacy-degree categories;
an external national higher-vocational reference completes the broad types. \\
STATENT & Establishment and enterprise evidence in the supplied extract &
School coordinates, NOGA compatibility, institution/campus grouping and
employment/FTE weights. Employment is a location proxy, not measured student capacity. \\
GWR residence candidates & External reviewed reference file & Descriptive
accommodation proximity for recovered residents. Proximity does not establish
student status or university affiliation. \\
\bottomrule
\end{longtable}

SS weights are used for fitting distributions and models. A pooled sum of three
annual SS weights is not a one-year population estimate: stock comparisons
must select a year or report an annual average. MZ person weights are used for
person-level reference distributions. Synthetic sample counts are divided by
\code{input_downsampling} when reporting full-population estimates; the result
is an expanded estimate, not an exact census count.

The supplied 2023 education codebook is also assumed to apply to 2021--2022;
the processed survey retains a flag for that assumption. SS usual commutes,
MZ diary-day visits and BFS institution enrolment have different statistical
scopes. Agreement between their totals does not prove that their frames or
weighting procedures are identical.

\section{Workflow and model order}

The education-relevant dependency order is:
\begin{enumerate}
\item Construct the STATPOP population and households, including the agreed
recovery and home-only rules.
\item Prepare SS labels and a separate student-training geography. Run the
existing employment/job and income models before predicting student participation.
\item Predict student status, then impute a joint study-group/study-country
outcome for existing students aged 15+.
\item Complete the intervening population models, optionally scale households
to projection controls, and sample the population.
\item Assign work and usual education locations before MZ donor matching.
Education assignment is independent of the future activity-chain donor.
\item Match MZ diaries using demographic/status compatibility and commute
information. Subsequent models, including public-transport subscription
ownership, consume the matched/enriched population.
\item Build daily trips and activities. Education visits use the assigned
education anchor. Diagnose student stock, campus anchors, daily visitors and distances.
\end{enumerate}

There is one finite education anchor per sampled person classified as a
student. This is enforced at the education-location stage boundary. The
foreign/unknown cases described below may have proxy anchors; a finite anchor
alone does not mean that its location is an observed domestic study destination.

\section{Student participation and apprenticeship}
\subsection{Participation label and predictors}

The binary student model uses the SS \code{CURRACTIVITY_STUDENT} label,
retained as \code{is_student}. \code{ONGOINGEDUCATION}, retained as
\code{current_education}, describes the programme and is used by the subsequent
classification model. The code preserves discrepancies between participation
and programme reports for auditing rather than declaring everyone with an
education programme to be a student.

The active classifier is CatBoost, fitted with survey weights. Its predictors
are age, age squared, sex, employed status, apprentice status, municipality
type, district and canton. Apprentice status means SS job-position code 50.
Municipality ID is used in data preparation but is \emph{not} an active
classifier predictor. Status is sampled from predicted probabilities, which
are also retained for diagnostics. There is no municipal student quota or
national BFS stock correction in this step.

For SS student observations, training geography uses the domestic school
\emph{departure} location when available. Otherwise it uses registered
residence. Explicit foreign departures remain an outside-Switzerland category.
Nonstudents retain residence geography. All relevant administrative features
are changed together on a separate training copy; other SS consumers retain
their original residence fields. Prediction uses the synthetic person's
current home geography. This departure/residence distinction remains a
limitation where students live elsewhere during term time.

The current student stage sets \code{is_student = 1} below age 15 after
prediction. This includes children below MZ interview age; it is a pipeline
convention, not evidence that every infant attends school. The flag, an anchor
and a diary education visit must therefore be counted separately.

\subsection{Partial repair of vocational records}

The former blanket rule made every SS vocational-secondary record an employed
apprentice. The current rule only considers age-15+ students in programme 6
who report nonemployment/inactivity and the corresponding job codes 60 or 70.
Explicit apprentices and other reported jobs are preserved.

Within each survey year, eligible records are shuffled reproducibly and the
prefix closest to the requested weighted repair mass is selected. Selected
records receive employed status and job-position code 50. Student flags and
survey weights do not change. The default repaired fraction is 0.65 of the
\emph{eligible weighted pool}, not 65\% of all vocational students.

This remains a scenario assumption motivated by the earlier stock audit, not
an observed apprenticeship label or an exact BFS calibration. Employment and
student stages export aggregate apprenticeship diagnostics separately for
survey years, private households and recovered residents. Apprenticeship also
enters programme imputation, reducing confusion between apprentices and other
employed tertiary students without imposing a deterministic programme rule.

\section{Study classification and age handling}

For students aged 15+, the programme model draws a joint programme/country
outcome from weighted SS profiles. Country status is domestic, abroad or
unknown. An explicit foreign school country takes precedence over a domestic
municipality field. Missing programme and country evidence remains visible.

The conditioning hierarchy is age band, employment, apprenticeship, sex and
canton, in that order. Each supported finer profile is blended with its
parent using
\[
 n_{\mathrm{eff}}=\frac{(\sum_j w_j)^2}{\sum_j w_j^2},\qquad
 P_g=\frac{n_{\mathrm{eff},g}}{n_{\mathrm{eff},g}+s}P_g^{\rm empirical}
       +\frac{s}{n_{\mathrm{eff},g}+s}P_{\rm parent},
\]
where \(s=20\) by default. The broadest age profile is empirical. Unseen
combinations back off to supported profiles; an unsupported age band uses the
nearest supported age band. Parent profiles include child observations, so this
is smoothing, not an independent Bayesian prior.

\begin{longtable}{L{.19\textwidth}L{.32\textwidth}L{.41\textwidth}}
\toprule Population group & Classification & Location handling \\
\midrule\endhead
Under 15 & Explicit below-SS-coverage marker & Existing age/school pools and
independent weighted MZ distance draws; no SS programme prediction. \\
Age 15+, SS codes 2--8 & Retain compulsory, bridging and secondary programme
detail & Separate pooled secondary OD model, SS age/programme distance targets
and compatible school selection. \\
Age 15+, SS codes 9--13 & One pooled tertiary outcome before fitting & Shared
UH, FH/PH and higher-vocational teaching pool; pooled tertiary OD and campus balancing. \\
Tertiary under 24 & Same pooled tertiary classification & Additional nationally
fitted distance preference, by default, before joint balancing. \\
Tertiary age 24+ & Same pooled tertiary classification & Original pooled seed
probabilities; shared campus balancing can still alter final destinations. \\
Foreign, unknown or unresolved study & Keep explicit status & Labelled
compatible or age-based proxy locations where regular domestic assignment does
not apply. These do not enter observed domestic OD evidence. \\
Nonstudents & No regular student anchor & Donor education visits are handled
as secondary \code{other} activities, with original-purpose provenance. \\
\bottomrule
\end{longtable}

The tertiary group combines higher professional examinations, higher vocational
colleges, bachelor, master and doctorate/habilitation. The synthetic population
does not receive separate tertiary degree labels: \code{current_education}
is missing for the pooled tertiary outcome and \code{study_level} is
\code{tertiary}. Detailed survey and BFS degree categories remain useful audit
references, not modelled degree stocks. FH/PH campuses are included in the
regular tertiary choice set.

Different stages use different age bands. Programme prediction uses
15--17, 18--19, 20--24, 25--29, 30--39, 40--49, 50--64 and 65+.
Secondary distance fitting uses 15--17, 18--23, 24--29, 30--39, 40--49,
50--59 and 60+. The tertiary tilt has only the under-24/24+ split.
Distance-analysis bands are not automatically the matching bands.

\section{Education origins and smoothed OD flows}

There are two active broad flow models: secondary (SS codes 2--8) and tertiary
(codes 9--13). Both use positive-weight student observations aged 15+ from
the pooled SS years. Origins prefer school-commute departure, with residence
fallback only when departure is unresolved. Explicit foreign origins and
foreign study destinations are excluded from domestic flows. Municipality
and quarter information is retained at the available resolution.

For an origin with empirical destination profile \(P_i\), the smoothed seed is
\[
 P_i^{\rm smooth}=\alpha_i P_i+(1-\alpha_i)R_i,
 \qquad \alpha_i=\frac{n_{\rm eff,i}}{n_{\rm eff,i}+s_{\rm OD}}.
\]
The default \(s_{\rm OD}\) is 100. A geographic prior \(R_i\) uses up to
50 observed neighbouring municipalities, a Gaussian bandwidth of 20 km and
capped effective-sample contributions so that one large city does not dominate.
For quarters, parent-municipality evidence excludes the quarter's own records
before regional borrowing. A zone with no observed respondents can therefore
still have a modelled destination distribution if the synthetic population
contains students there. Borrowed evidence is recorded as such.

A weak distance penalty applies to uncertain origin profiles:
\[
 P_{ij}^{\rm base}\ \propto\ P_{ij}^{\rm smooth}
 \exp\{-\lambda(1-\alpha_i)d_{ij}\},
 \qquad \lambda=0.005\ {\rm km}^{-1}.
\]
Distances here are zone-centroid straight-line distances. The implementation
centres exponents for numerical stability. Geographic borrowing can introduce
routes absent from an origin's own sample; subsequent distance weighting and
balancing preserve zeros in the resulting smoothed seed. Preserving support
does not mean preserving every zero in the raw survey count table.

\section{Child and secondary-school assignment}
\subsection{Children and baseline proxies}

The legacy location pools use age intervals \(a\leq6\), \(6<a\leq12\),
\(12<a\leq16\), and \(a>16\), associated with kindergarten, primary,
secondary and tertiary establishments. The initial nearest-site candidate
counts are 1, 2, 7 and 20 respectively. Classified age-15+ domestic students
are subsequently handled by the secondary/tertiary methods below.

Independent weighted MZ education-distance draws use these same age intervals,
with the configured diary-day selection. They do not require the person's
future donor. The youngest pool may use age-six evidence if present; MZ has
no observations below six. An empty reference pool receives a zero-distance
reference and an explicit nearest-school coverage-fallback label. This is
particularly relevant to the kindergarten pool, whose candidate count is one.

MZ diary distances use the last observed home before the selected education
activity, otherwise the first home afterwards, and finally the household home.
Secondary-home diary activities are eligible. The main education activity is
selected by duration, with origin-only education handled when necessary.

\subsection{Secondary programmes and distance targets}

The stage \code{secondary_distances} fits weighted empirical SS distance CDFs
for each age/programme cell. It uses reported school-departure geography and
nonnegative \code{DISTTOSCHOOL_AIR}, interpreted as kilometres by default.
Residence-fallback OD observations do not supply these reported-departure
distance references. Codes 2 and 3 share one programme pool. Below effective
sample size 50, fitting backs off to all ages of the same programme. Weak
programme-wide support is labelled; absent evidence gives a missing target
and nearest-compatible-school fallback.

Within each recipient age/programme group, midpoint quantiles of the fitted
CDF supply target distances. Independent MZ references determine rank, with
seeded handling of ties/missing references. Thus SS determines the distance
distribution; the MZ draw is not the literal target distance.

\begin{center}\small
\begin{tabular}{lll}
\toprule SS code & Programme & Compatible establishment NOGA \\
\midrule
2--3 & Compulsory & 853101, 852002, 852003 \\
4 & Bridging & 853101, 853102, 853103, 853200 \\
5 & Specialised secondary & 853102, 853103 \\
6 & Vocational secondary & 853200 \\
7 & General secondary & 853102 \\
8 & Vocational/specialised maturity & 853200, 853102, 853103 \\
\bottomrule
\end{tabular}\end{center}

NOGA compatibility is an explicit approximation, not a verified programme
catalogue. In particular, code 5 retains both academic school types because
the narrower establishment code alone has insufficient coverage.

The final secondary destination zone is drawn from the compatible pooled OD
row after multiplying by
\[
 \exp\left[-\frac12\left(\frac{e_{ij}}{\sigma_i}\right)^2\right],
 \quad \sigma_i=\max(1000\ {\rm m},0.10t_i),
\]
where \(t_i\) is the SS target distance and \(e_{ij}\) is the smallest absolute
target mismatch among compatible schools in zone \(j\), measured from the
actual synthetic home. The zone receives no artificial bonus for duplicate
school records. Missing targets leave compatible OD probabilities unconditioned.

Inside the sampled zone, up to 20 schools with the smallest target-distance
mismatch are considered. A Gaussian preference with tolerance
\(\max(200\ {\rm m},0.10t_i)\) is combined with clipped employee weights.
Unsupported OD choices retain the preceding compatible distance-based
fallback. Scarce schools can prevent exact reproduction of targets. There
are no secondary school-capacity or destination-total constraints.

\section{Tertiary institutions, campuses and balancing}
\subsection{Establishment registry and BFS reference}

STATENT establishment and enterprise NOGA identify UH (854201), FH/PH
(854202) and higher-vocational (854203) institutions. Establishments are
grouped within an enterprise using DBSCAN with default 500 m neighbourhood
distance and one-point minimum clusters. This is a connectivity rule, not a
maximum campus diameter; chained sites can span more than 500 m.

Regular pooled assignment uses eligible teaching establishments, excluding
type conflicts and applying reviewed site rules. Research-only sites remain
outside this undifferentiated teaching pool. Detailed doctoral eligibility
helpers remain available for analysis but are not a separate active doctoral
assignment model.

External institution links associate BFS totals with STATENT enterprises.
Both proposed and verified link statuses can currently supply targets; their
evidence quality should be inspected. Within a linked institution, eligible
teaching-site FTE splits the total between campuses:
\[
 E_c=E_u\frac{F_c}{\sum_{k\in u}F_k}.
\]
Here \(E_u\) is institution enrolment and \(F_c\) is eligible campus FTE.
\(E_c\) is an estimated campus allocation, not an independently published
campus enrolment count.

The pooled annual reference includes disjoint university programmes, including
continuing/other and legacy diplomas, plus the supplied higher-vocational
reference. The documented 2023/24 input has 169,249 UH, 107,196 FH/PH and
60,854 higher-vocational enrolments (337,299 combined). These are reference
file values, not a claim of complete coverage of every tertiary training
activity or a Swiss-resident census. Changing the input/year requires rechecking them.

\subsection{Targets and assignment}

For \(N\) modelled domestic tertiary students, the initial campus target is
\[
 T_c=N\frac{E_c}{\sum_u E_u}.
\]
Named campuses and residual institution-type groups form disjoint margins.
Unlinked enrolment stays in the residual UH/FH/PH/higher-vocational groups;
it is not redistributed exclusively to named institutions.

Balancing operates on the cohort with usable compatible OD support, using
that supported cohort's size for its requested margins. Unsupported people
retain labelled proxy/fallback anchors outside those fitted margins. Iterative
proportional fitting adjusts origin-to-campus/type flows. Controlled rounding
preserves each supported origin-row total and the fitted campus/type margins
within the integer rounding constraints. Infeasible targets remain visible
as residuals; the algorithm does not invent routes to satisfy them.

With the default young-distance feature, rows are origin-by-age-group, and
the young seed is additionally multiplied by
\[
 \exp(-\beta d_{ij}),\qquad a_i<24.
\]
One national \(\beta\) is fitted by conditional destination likelihood from
reported domestic school departures, with a numerical bound
\(0\leq\beta\leq0.05\ {\rm km}^{-1}\). Reported numeric SS distance is not
required: the fit uses observed OD-zone centroids. The age threshold and bound
are code constants. Older students have no added tilt, although balancing
their shared campus margins can change final allocations. The October audit
fit was approximately 0.00887 per kilometre; this is a fitted result, not a
configuration default.

After selecting a zone and calibration group, named groups identify their
campus; residual groups can contain several campuses, selected by eligible
FTE before individual buildings. This avoids a large cluster of nearby
buildings excluding another campus. Within a chosen campus/zone, the existing
employee and reference-distance preference selects among up to 20 nearest
buildings. FTE controls campus attractiveness; headcount contributes to the
within-campus preference. Neither is a hard building capacity.

\section{Matching and daily attendance}

Work and education anchors exist before matching. For donors and recipients,
the combined matching predictor is the longer valid work/education distance:
\[
 d_{\rm match}=\max(d_{\rm work},d_{\rm education}).
\]
If only one is known, that one is used; if neither is known, the sentinel is
\(-1\), while zero is a valid commute. On the MZ side, school distance uses
the reported usual school address first, then an observed diary education
location. Remaining missing student distances are drawn with person weights
from observed domestic peers: same age band and employment first, then age
band, then all students. Imputed distances never become donor references for
further imputation and never manufacture education visits.

Current ordinary matching uses broad donor-pool splits under 15, 15--23 and
24+. Status and age-class variables further constrain matching; the exact
matching classes are separate from the distance-report bands. For ages 15--23,
employment and combined employment/student status are mandatory, while
commute classes and canton are among the relaxable criteria. The algorithm
can relax detailed matching requirements and reduce the requested donor count;
it does not guarantee a same-canton donor or an exact distance match.

Daily attendance follows the selected diary. A student can have no education
visit because the observed day includes work, other activities or staying
home. A workday selection includes weekdays during school holidays. There is
no additional fixed attendance-rate draw, national attendance multiplier or
daily BFS campus quota in the active assignment.

For nonstudent donors, reported education-purpose visits are recoded to
\code{other} in the secondary-trip processing, retaining original-purpose
provenance. Synthetic trip processing also applies the recoding using the
recipient's final nonstudent status; raw MZ diaries remain available unchanged
for matching and survey audits. This allows secondary activity assignment to handle possible
courses and continuing education without assigning every nonstudent a regular
school. It does not establish the actual nature of those courses.

\section{Recovered residents, scaling and residence limitations}

Geocoded people missing from the household table can be retained in technical
single-person households. Recovery requires the existing coordinate-quality
checks; recovered children under 16 are excluded. Younger recovered people
use the same ordinary employment, student and programme models as comparable
private-household residents. Dorm proximity is descriptive. There are no
forced dorm-student labels, donor swaps, reserved campus places or subtraction
of a fixed dorm group from municipal student quotas.

With residence classification enabled, recovered people aged 39+ are assigned
the agreed home-only scenario and bypass the corresponding behavioural
models/diagnostics. This age cutoff is a modelling assumption, not an observed
absence of mobility. The older home-only rule is independent of dorm distance.

Optional projection scaling uses year-end population controls from the
2025 reference vintage and compatible private-household controls. Technical
recovered households are excluded from the private-household margin but their
people remain in population controls and household replication. Scaling changes
the cohort composition and hence can change student stocks without changing
the student model itself. For direct 2023 comparisons, use scaling disabled.

Raw STATPOP secondary-residence fields and linked secondary-coordinate records
were identified during the latest audit, but \textbf{no term-time relocation
has been implemented}. The retained home remains the current registered-home
representation. School-departure geography in SS training/ODs therefore does
not automatically move synthetic students to term-time homes. Multiple raw
residence records must never be counted as additional people without resolving
identity and residence scope. This is a documented future investigation, not
an active configuration option.

\section{Configuration reference}
\label{sec:config}

\textbf{Conventions.} Defaults are stage defaults when the relevant stage is
configured. Required general inputs such as \code{data_path},
\code{analysis_path}, \code{random_seed} and \code{input_downsampling} must
come from the surrounding configuration. Numeric parameters must remain in
their documented valid ranges. Booleans should be YAML \code{true}/\code{false}.
Changing defaults is a sensitivity experiment, not automatically an improvement.

\subsection{Participation, programme and distance settings}
\small
\begin{longtable}{L{.34\textwidth}L{.13\textwidth}L{.45\textwidth}}
\toprule Parameter & Default & Alternatives and effect \\
\midrule\endhead
\code{vocational_apprentice_repair_fraction} & 0.65 & Range [0,1]. Zero retains
raw labels; one repairs all eligible ambiguous records. Even one does not
restore the old blanket overwrite of every vocational record. Does not change student flags. \\
\code{education_level_prior_strength} & 20.0 & Nonnegative effective-observation
strength. Larger values shrink finer programme/country profiles more strongly;
zero removes blending for supported cells. Unsupported profiles still back off. \\
\code{secondary_education_distance_minimum_effective_n} & 50 & Positive threshold
for age/programme distance cells. Larger values pool more cells across ages;
smaller values retain noisier local age/programme distributions. \\
\code{education_distance_ss_unit} & \code{km} & \code{m} is the alternative.
Controls SS reported-distance conversion in secondary targets and distance
analysis; it does not change centroid-distance units. Set from input metadata,
not to tune results. \\
\bottomrule
\end{longtable}
\normalsize

\subsection{OD smoothing and regularisation}
These settings are independent for secondary and tertiary flows. Values shown
are identical defaults, but each key controls only its corresponding matrix.

\small
\begin{longtable}{L{.34\textwidth}L{.13\textwidth}L{.45\textwidth}}
\toprule Parameter & Default & Alternatives and effect \\
\midrule\endhead
\code{secondary_od_prior_strength} & 100.0 & Nonnegative shrinkage strength.
Higher values favour neighbouring/parent evidence. Zero uses local empirical
profiles where observed, but still needs borrowed evidence for empty origins. \\
\code{tertiary_od_prior_strength} & 100.0 & Same interpretation for tertiary flows. \\
\code{secondary_od_neighbors} & 50 & Positive integer maximum number of nearby
observed municipalities in geographic priors. More broadens the donor set;
fewer makes it more local and potentially less stable. \\
\code{tertiary_od_neighbors} & 50 & Same interpretation for tertiary flows. \\
\code{secondary_od_bandwidth} & 20000.0 m & Positive Gaussian distance scale
for neighbouring origins. Larger values spread influence more widely;
smaller values concentrate it locally. Not a hard search radius. \\
\code{tertiary_od_bandwidth} & 20000.0 m & Same interpretation for tertiary flows. \\
\code{secondary_od_distance_decay} & 0.005/km & Nonnegative weak distance penalty,
scaled by origin uncertainty. Zero disables it; larger values favour nearby
supported destinations more strongly. \\
\code{tertiary_od_distance_decay} & 0.005/km & Same base penalty for tertiary
flows, separate from the fitted under-24 preference. \\
\bottomrule
\end{longtable}
\normalsize

\subsection{Campus assignment and enrolment inputs}
\small
\begin{longtable}{L{.34\textwidth}L{.13\textwidth}L{.45\textwidth}}
\toprule Parameter & Default & Alternatives and effect \\
\midrule\endhead
\code{education_tertiary_balance} & \code{true} & False uses compatible pooled
tertiary OD without campus/type margin balancing and also disables the young
distance preference. It does not remove tertiary anchors. \\
\code{education_tertiary_age_distance} & \code{true} & Active only with balancing.
False retains pooled OD seeds without the added under-24 tilt, while keeping
campus balancing. It does not disable the base OD distance penalty. \\
\code{education_campus_cluster_distance} & 500.0 m & Positive within-enterprise
DBSCAN neighbourhood scale. Smaller values split campuses; larger values merge
nearby site clusters. Changing it changes campus IDs/groups and requires reviewing maps and targets. \\
\code{education_enrolment_year} & \code{2023/24} & Choose another academic-year
string only if matching university and higher-vocational input rows are available.
Changing the year does not download data or rescale the population year. \\
\code{education_enrolment_path} & \code{null} & Null resolves to
\code{<data_path>/education_enrolment}. An explicit directory selects another
set of saved BFS JSON snapshots. The CSV overrides below have independent defaults. \\
\code{education_institution_links} & Path A & Explicit CSV path changes
BFS-to-STATENT institution mapping. Missing/unreviewed links affect named margins
and residual groups; do not duplicate one institution's total across enterprises. \\
\code{education_site_eligibility} & Path B & Explicit CSV path selects reviewed
site/programme eligibility decisions. Changes can affect teaching pools and
campus FTE weights. Rules need supported IDs and evidence. \\
\code{education_higher_vocational_reference} & Path C & Explicit CSV path selects
the national higher-vocational reference. Changes alter its share in pooled
tertiary targets. It is not a measured campus-capacity file. \\
\code{education_reference_sites} & Path D & Reference-site CSV for institution
matching and campus audits. It is not an automatic home-relocation or student
affiliation rule. \\
\bottomrule
\end{longtable}
\normalsize

Paths A--D are under \code{<data_path>/education_enrolment/}:
\begin{description}
\item[A] \code{education_institution_links.csv}
\item[B] \code{education_site_eligibility.csv}
\item[C] \code{education_higher_vocational_enrolment.csv}
\item[D] \code{education_reference_sites.csv}
\end{description}
These files remain external data inputs, not committed reference CSVs.
Use explicit valid filenames to override these four path defaults; unlike
\code{education_enrolment_path}, null is not their documented fallback value.

\subsection{Related population, matching and analysis controls}
\small
\begin{longtable}{L{.34\textwidth}L{.13\textwidth}L{.45\textwidth}}
\toprule Parameter & Default & Alternatives and effect \\
\midrule\endhead
\code{retain_geocoded_missing_households} & \code{true} & False excludes missing-
household records instead of retaining eligible geocoded residents as technical
singletons. Changes population coverage and potentially education stocks. \\
\code{non_private_residences_enabled} & \code{true} & False disables GWR proximity
classification and its older-resident home-only flags. Recovery remains governed
by the preceding setting; retained residents then use ordinary model handling. \\
\code{non_private_residence_radius} & 250.0 m & Nonnegative descriptive proximity
threshold. Larger values label more recovered younger people as dorm candidates.
This flag does not force them to become students or select their campus. \\
\code{non_private_home_only_minimum_age} & 39 & Must exceed 18. Recovered people
at/above it are home-only; younger recovered people use ordinary models. Also
sets the exclusive upper age for the dorm-candidate flag. \\
\code{gwr_residence_candidates} & Path E & Defaults to
\code{<data_path>/gwr/residence_candidates.csv}; override with a compatible
reviewed residence-candidate file. \\
\code{only_permanent_residents} & \code{true} & True admits residence type 1;
false admits types 1 and 2 at the first filter. Both branches subsequently
require population type 1. False therefore does not by itself retain all
secondary-residence records or implement term-time homes. \\
\code{specific_day_scenario} & \code{workday} & \code{weekend} or capitalised
weekday names Monday--Sunday select other diary pools. Affects matching and
MZ education references/analysis, not BFS annual enrolment. \\
\code{matching_minimum_observations} & 10 & Requested donor support before
progressive matching relaxation/reduction. Larger values can cause more relaxation;
not a guaranteed minimum in the final most-specific cell. \\
\code{matching_enable_diagnostics} & \code{false} & True enables additional
matching feature-distribution logs. Does not impose additional matching targets. \\
\code{education_stock_survey_year} & 2023 & Selects the year for the stock audit.
It does not switch production training from pooled SS to a single year. \\
\code{enable_scaling} & \code{false} & True fits and replicates households using
projection controls. False preserves the unscaled source population before sampling. \\
\code{scaling_year} & 2050 & Existing code default, relevant only if scaling is
enabled. Explicitly set 2025 for 31 December 2025 targets from the 2025 reference
vintage. The target year is not implicitly 2025. \\
\code{input_downsampling} & Required & Sample fraction used by population
sampling and expansion of analysis counts. Set consistently; do not compare
raw sample counts directly with full BFS stocks. \\
\code{random_seed} & Required & Used by programme draws, OD/location sampling,
apprentice repair and other stages. Does not override every legacy model seed;
see Section~\ref{sec:constants}. \\
\bottomrule
\end{longtable}
\normalsize

Scaling inputs are unchanged official BFS CSV exports under
\code{<data_path>/projections/raw/}. The stage
\code{data.statpop.projections.raw} decodes and combines these files before the
population and household stages build target-year controls. The previously
prepared \code{population_year_end_2025.csv} and \code{households_2025.csv} are
no longer required. Download and sharing instructions are in
\code{doc/data/statpop/projections.md}. With scaling disabled, these raw inputs
are not required.

\subsection{Important choices that are not YAML switches}
\label{sec:constants}

The current code fixes departure-first student-training geography, pooled
tertiary classification, and the pre-matching education-anchor order. No active
switch restores forced dorm students, municipal dorm subtraction, dorm swaps
or separate bachelor/master/doctoral assignment.

Other code constants include the under-24 threshold and coefficient bound,
secondary zone/school tolerance floors (1000/200 m), the 10\% relative tolerance,
and the school candidate counts. The student CatBoost settings are 3200
iterations, learning rate 0.04, depth 10, L2 regularisation 6, Bernoulli
bootstrap and subsample 0.8, with training seed 42. The student draw uses the
legacy base seed 789 plus chunk offset, with chunks of 50,000. These are not
controlled by the general \code{random_seed} setting. Changing the general
seed should not be described as refitting every model with a new seed.

\subsection{Example education configuration}

Merge the following into the existing YAML \code{config:} mapping. This is
an education-focused fragment, not a complete runnable pipeline file.
The explicit 2025 target below overrides the legacy 2050 default but has no
effect while scaling is disabled.

\begin{lstlisting}
config:
  vocational_apprentice_repair_fraction: 0.65
  education_level_prior_strength: 20.0

  secondary_od_prior_strength: 100.0
  secondary_od_neighbors: 50
  secondary_od_bandwidth: 20000.0
  secondary_od_distance_decay: 0.005
  secondary_education_distance_minimum_effective_n: 50
  education_distance_ss_unit: km

  tertiary_od_prior_strength: 100.0
  tertiary_od_neighbors: 50
  tertiary_od_bandwidth: 20000.0
  tertiary_od_distance_decay: 0.005
  education_tertiary_balance: true
  education_tertiary_age_distance: true

  education_campus_cluster_distance: 500.0
  education_enrolment_year: "2023/24"
  education_enrolment_path: null
  # Reviewed CSV paths default independently under data_path.

  retain_geocoded_missing_households: true
  non_private_residences_enabled: true
  non_private_residence_radius: 250.0
  non_private_home_only_minimum_age: 39

  specific_day_scenario: workday
  education_stock_survey_year: 2023
  enable_scaling: false
  scaling_year: 2025
\end{lstlisting}

\section{Validation and interpretation}

\subsection{Checks that answer different questions}
\begin{enumerate}
\item \textbf{Population stock:} count all students regardless of diary visits;
compare age-15+ totals and employment/student overlap with SS 2023. Separate
private, recovered and home-only groups; retain abroad/unknown study status.
\item \textbf{Programme prediction:} examine effective sample sizes, backoff
methods, apprenticeship conditioning and held-out outcomes. Detailed degree
counts cannot be validated as synthetic outcomes when tertiary is pooled.
\item \textbf{Annual anchors:} verify one finite location per student, compatible
schools, origin support, campus/type margins and unsupported-proxy shares.
\item \textbf{Daily attendance:} count distinct people with an education activity,
not repeated visits. Compare modelled diary attendance with weighted MZ,
standardising for age/status composition before attributing a difference to placement.
\item \textbf{Distance:} compare means, medians, tails and empirical CDFs.
MZ uses diary-home distance; SS uses reported departure-based distance; the
population uses assigned home and actual education visitors. A separate common
municipality-centroid comparison removes differences in coordinate precision.
\end{enumerate}

The young-tertiary distance audit rebuilds both base OD and coefficient on two
SS years and tests the omitted year, rotating 2021--2023. It compares baseline,
half and full fitted preference using weighted log loss, Brier score and
zero-support share. The documented trial favoured the full young preference
in all three held-out years, with coefficients about 0.00885--0.00940/km.
This supports a modest national age effect, not precise local prediction or
resolution of term-time residence.

Regression checks cover origin/age counts, structural support, integer campus
margins, study/status preservation, geographic ambiguity, apprenticeship
repair and the real Synpp configuration/execute interface. The commit review
ran 49 targeted tests plus 17 subtests across the recent changes, including
related freight and simulation fixes. Passing implementation tests is not a
statistical guarantee of every campus's attendance.

\subsection{Analysis stages and outputs}
Outputs are generally placed under \code{<analysis_path>/<analysis_name>}.
The production run must include the relevant downstream activities to report
actual attendance; a frozen-diary replay only isolates placement effects.

\begin{longtable}{L{.44\textwidth}L{.49\textwidth}}
\toprule Stage & Main use \\
\midrule\endhead
\code{analysis.tertiary_education_assignment} & Annual anchors, campus counts,
origin flows, daily visitors and balancing diagnostics. Inspect
\code{calibrated_group_comparison.csv} and
\code{student_anchor_and_attendance.csv}. \\
\code{analysis.education_student_stock} & Stock by residence cohort, student/
employment overlap, SS-year comparisons and BFS reference coverage. Audit only;
does not impose stock corrections. \\
\code{analysis.education_distances} & Weighted native and common-centroid
distance summaries/CDFs by age and student scope, with coverage diagnostics. \\
\code{analysis.secondary_education_assignment} & Secondary compatibility,
targets, realised distances and fallback methods. Some donor-based columns
are legacy counterfactual diagnostics, not current assignment objectives. \\
\code{analysis.tertiary_education_age_distance} & Held-out-year validation of the
national young-distance preference. \\
\code{analysis.education_classification} & Smoothed programme/country prediction
and support; does not retrain the binary student classifier. \\
\code{analysis.education_institutions} and
\code{analysis.education_campus_weights} & Registry clustering, institution
coverage, FTE allocation and reviewed site eligibility. \\
\bottomrule
\end{longtable}

An ``expected daily'' number displayed on a comparison map is an analytical
benchmark obtained by applying an attendance assumption/rate to enrolment.
It is not the active campus-balancing target. A ``cohort-normalised target''
instead describes the model's annual student-anchor margin. Neither should be
labelled a directly published campus attendance count.

\subsection{Remaining limitations and current stopping point}

The main uncertainties are differing source years and frames, finite local
survey support, the apprenticeship repair assumption, FTE-based campus shares,
secondary NOGA compatibility, missing-distance imputation and registered versus
term-time homes. Pooling tertiary prevents sparse degree-specific OD matrices,
but removes degree-specific placement, including a separate research-site
model for doctoral students. Continuing education also has incomplete alignment
between participation labels, BFS categories and MZ purposes.

The agreed baseline keeps the current implementation. Secondary-residence
relocation and additional local tuning are deferred. Further changes should
be justified by a measured residual with adequate support, rather than by
forcing every campus or small canton/age cell to match one reference exactly.

\appendix
\section{Implementation guide}
Paths below are relative to the repository root. They identify the active
implementation, rather than requiring external documentation to compile the report.
\begin{description}
\raggedright
\item[Survey preparation] \code{data/structural_survey/raw.py},
\code{structural_survey.py}, \code{student_training.py},
\code{education_origins.py} and \code{apprenticeship.py} in the same directory.
\item[Population models] \code{synthesis/population/models/students.py},
\code{education.py}, \code{employment.py} and \code{non_private.py}.
\item[OD evidence] \code{data/od/education_matrix.py},
\code{tertiary_education_matrix.py}, \code{tertiary_education_smoothing.py}
and \code{tertiary_education_age.py}.
\item[Institutions and references] \code{data/statent/education_institutions.py},
\code{education_enrolment.py} and \code{education_campus_weights.py}.
\item[Location assignment] \code{synthesis/population/spatial/primary/education/}
contains \code{locations.py}, \code{secondary_distances.py},
\code{secondary.py}, \code{targets.py}, \code{balancing.py} and \code{tertiary.py}.
\item[Matching and diary purposes] \code{synthesis/population/matching/commute.py},
\code{matched_v1.py}, \code{data/microcensus/commute.py} and
\code{data/microcensus/secondary_trips.py}.
\item[Recovery and projections] \code{data/statpop/statpop.py},
\code{data/statpop/households.py}, \code{data/statpop/projections/} and
\code{synthesis/population/scaled.py}.
\item[Companion documentation] The \code{analysis/education_*.md},
\code{analysis/tertiary_education_*.md} and secondary-assignment notes provide
stage-level details. Some include historical trials; active behaviour is
resolved against the executable stage entry points in this report.
\end{description}

\section{Online sources for external reference files A--E}

The five CSV filenames in Section~\ref{sec:config} are project-prepared inputs,
not official files available online under those names. Official sources supply
enrolment tables, campus addresses and building records; our crosswalks, review
decisions, aggregation and filtering turn them into the required CSV schemas.
Preserve the reviewed snapshots to reproduce the current run. The normal pipeline
reads local inputs and does not rebuild them from live websites.

\begin{description}
\raggedright
\item[A: Institution links.] \code{education_institution_links.csv} is our
reviewed crosswalk from named BFS institutions to anonymous STATENT enterprise
identifiers. There is no direct public download for this crosswalk. Evidence
includes \href{https://www.epfl.ch/campus/visitors/coming-to-epfl/}{EPFL campus addresses},
\href{https://www.zhaw.ch/en/about-us/visitor-information/location-plans}{ZHAW location plans},
\href{https://ethz.ch/en/campus/erreichen/hoenggerberg.html}{ETH Hoenggerberg}
and \href{https://students.uzh.ch/de/dates/sites/irchel.html}{UZH Irchel}.
Public addresses support the spatial interpretation but do not independently
verify an anonymised enterprise identity. Use the existing reviewed mapping
with the corresponding STATENT extract.

\item[B: Site eligibility.] \code{education_site_eligibility.csv} records our
modelling decisions, with per-row evidence and notes. It is not an official
campus-eligibility dataset. Examples of source evidence are
\href{https://www.epfl.ch/education/bachelor/enhanced-learning-programs/summer-in-the-lab-internship-program/}{EPFL's laboratory internship programme}
and \href{https://csf.ethz.ch/monte-verita.html}{ETH's Monte Verita conference venue}.
These sources describe activities/use; the decision about regular student
assignment remains an explicit project assumption.

\item[C: Higher-vocational enrolment.]
\code{education_higher_vocational_enrolment.csv} is a compact transformation of
\href{https://www.pxweb.bfs.admin.ch/pxweb/de/px-x-1502030000_112/-/px-x-1502030000_112.px/}{BFS table px-x-1502030000\_112}.
Select academic year 2023/24 and aggregate across fields, sexes and nationality
categories without counting total rows alongside their components. Retain the
education-type distinction. Our snapshot contains 35,982 higher-vocational
college students, 21,960 in the examination/preparation category and 2,912 in
non-BBG-regulated higher vocational education. The file preserves the source
URL and incomplete preparation-course coverage note. A raw BFS export needs
conversion to the pipeline's column names and three study-level labels.

\item[D: Campus reference points.] \code{education_reference_sites.csv} was
assembled from official campus maps/addresses, including the institution
pages above. The \href{https://plan.epfl.ch/}{EPFL map} and
\href{https://students.uzh.ch/de/dates/sites/irchel.html}{UZH Irchel map/address}
are examples. Address coordinates were obtained using swisstopo's location
search service; see the \href{https://docs.geo.admin.ch/}{geo.admin.ch API documentation}.
The CSV retains source URLs, coordinate sources, notes and review dates. These
are campus anchors, not campus polygons or an official combined national file.

\item[E: Residence candidates.] \code{gwr/residence_candidates.csv} is our
filtered and enriched building snapshot. Download raw GWR data from
\href{https://public.madd.bfs.admin.ch/ch.zip}{the nationwide GWR ZIP}; see
\href{https://www.gwr.admin.ch/de/data/supply/public_content.html}{official access and field documentation}.
Operator evidence comes from \href{https://www.fmel.ch/nos-maisons}{FMEL's residence list}.
Our 13 October 2026 snapshot retains communal-residence class 1130, additional
explicit student-housing names and reviewed FMEL address matches, deduplicated
by EGID. It contains 4,774 buildings. Raw GWR alone is neither this prepared
CSV nor a complete student-dormitory inventory; communal accommodation includes
other uses. Live exports change over time, so re-downloading will not necessarily
reproduce the saved snapshot.
\end{description}

\section{Reproducible maintenance of files A--D}

The standalone command \code{python -m analysis.education_reference_update}
provides two operations: \code{prepare} creates a new candidate reference
directory; \code{review} checks it against an exported STATENT registry. Neither
operation changes the active configuration or runs the population pipeline.
Detailed commands, file provenance, conversion rules and review procedures are
documented in \code{doc/analysis/education_reference_update.md}.

Preparation copies A, B and D unchanged, downloads the official UH, FH/PH and
higher-vocational cubes for explicitly requested years, and derives C. It
archives queries, metadata, responses, retrieval timestamps and hashes. An
offline mode reproduces the conversion from those saved responses. The C
converter selects only totals on the field/sex/nationality axes, maps education
types 311, 312+313 and 315 to the three reference groups, and reconciles their
sum against the published total. Missing cells, changed definitions or positive
unmapped historic categories require review rather than silent substitution.

Review checks missing or duplicate identifiers, type/eligibility consistency,
source evidence, C conservation, and proximity of D anchors to linked sites.
With a prior registry export, it also reports establishment, enterprise, NOGA
and coordinate changes. Its default 500\,m proximity threshold is only an audit
threshold, not a campus boundary or assignment parameter. Proposed identities
remain proposed; no automatic spatial remapping or promotion takes place.
The HTML report, issue tables and registry baseline remain in the external
data directory alongside the reference bundle.

A validation download on 21 October 2026 reproduced all three existing 2023/24
higher-vocational counts exactly (60,854 in total). Download and offline replay
produced identical A--D CSVs. The current cached STATENT registry produced no
structural errors, but retained five proposed-identity notices and flagged the
ZHAW Zurich Lagerstrasse anchor for spatial review. The first audit also records
the lack of a historical registry baseline. These checks do not certify the
anonymous institution identities. The existing production references and
assignment remain unchanged.

\end{document}
